\section{结果必须按方向、资源与指标拆开读}

模型训练完成后，最危险的做法是只报一个平均分。两百语言包含 into-English、out-of-English 和 non-English directions，也包含 high、low、very-low resource groups；平均值会让数量多或分数高的组掩盖真正薄弱的方向。课堂用三张表逐步训练这种分解阅读能力。

\subsection{先与前代多语系统比较}

本节从训练机制转入系统级结果，先建立历史基线。FLORES-101 对比表把 NLLB 与 M2M-100、DeepNet 等前代系统放在同一 evaluation setup。读表时应确认每个模型覆盖的语言数、参数/训练数据和 metric 是否一致，再看平均分变化。

\lecturefigure{slide-31-flores101-comparison.jpg}{NLLB 与前代多语系统在 FLORES-101 上的比较}{Stanford 课堂视频 00:32:06}

\begin{knowledgebox}{读图：先看列定义与缺失项，再看赢家}
表格分别列出 into-English、out-of-English、non-English 与平均结果；`-/-` 表示某些来源没有可直接比较结果。NLLB 的提升支持“数据、模型和训练联合改进”，却不能单凭一张表分离其中每个组件的独立贡献。
\end{knowledgebox}

\begin{warningbox}{模型代际比较要控制 evaluation protocol}
分词器、语言集合、checkpoint、解码设置和 metric implementation 都会改变分数。若协议不同，表中的数值只能作为课堂时间点的系统级参考，不应被重新排列成脱离上下文的永久 leaderboard。
\end{warningbox}

\subsection{FLORES-200：按资源层级拆分}

FLORES-200 表把 out-of-English 与 into-English 进一步分为 all、high、low、very-low，并单列 non-English。它让读者看到总体提升是否真正落到低资源尾部，而不是只来自高资源语言。

\lecturefigure{slide-32-flores200-results.jpg}{FLORES-200 按方向和资源层级报告的结果}{Stanford 课堂视频 00:33:09}

\begin{knowledgebox}{读图：先横向比较资源层级，再纵向比较方向}
同一行中 high 到 very-low 的落差表示资源条件影响；同一资源层中 into 与 out-of-English 的差异表示方向不对称；non-English 汇总依赖多语 transfer。平均值能概括系统，却不能替代这三种诊断。
\end{knowledgebox}

课堂结果同时出现 BLEU 与 chrF++。BLEU 主要比较 tokenized n-gram precision 并带长度惩罚；chrF++ 更依赖字符 n-gram，并加入 word n-gram，在形态丰富或 tokenization 不稳定的语言中常更稳健。二者都只是 reference-based automatic metrics。

\begin{knowledgebox}{术语消化：结果表中的指标}
\begin{tabularx}{\textwidth}{@{}L{0.2\textwidth}L{0.27\textwidth}X@{}}
\toprule
指标/分组 & 解决的问题 & 不能证明什么 \\
\midrule
BLEU / spBLEU & 以统一分词比较 n-gram 匹配 & 不直接代表语义、自然度或安全 \\
chrF++ & 用字符与词 n-gram 降低分词敏感性 & 仍依赖参考译文，不能覆盖所有合法表达 \\
Into-English & 多语言输入到英语 & 不能代表目标为低资源语言的生成质量 \\
Out-of-English & 英语输入到多语言 & 容易受目标端生成与数据稀缺影响 \\
Non-English & 两个非英语语言间翻译 & 常依赖 transfer，训练方向覆盖不均 \\
\bottomrule
\end{tabularx}
\end{knowledgebox}

\subsection{与商业系统比较的正确边界}

课堂再把低资源方向与 Google Translate 等外部系统放在一起。这个对比用于说明 NLLB 在研究覆盖和公开模型上的竞争力，但外部产品不断更新，支持语言、解码和内部数据也不可见，因此它不是稳定基准。

\lecturefigure{slide-33-external-system-comparison.jpg}{NLLB 与课堂时间点商业翻译系统的结果比较}{Stanford 课堂视频 00:33:27}

\begin{knowledgebox}{读图：产品对比首先是一张时间戳快照}
先确认哪些语言方向双方都支持，再比较低资源列；空缺不应视为零分。商业系统可能使用未公开数据和人工规则，NLLB 则强调开放研究资产。表格支持“在部分方向达到有竞争力质量”，不支持“所有场景全面优于产品”。
\end{knowledgebox}

\begin{importantbox}{44\% BLEU 是相对提升}
若前代分数为 $B_{\text{old}}$，新分数为 $B_{\text{new}}$，相对提升定义为
\[
r=\frac{B_{\text{new}}-B_{\text{old}}}{B_{\text{old}}}\times100\%.
\]
$r=44\%$ 不表示 BLEU 增加 44 个绝对点，也不表示每个语言对都提升 44\%。它是 NLLB 项目在指定比较集合与协议下报告的总体相对改进。
\end{importantbox}

\teachervoice{Fan 给出一个课堂经验判断：某些指标大约到 30 左右会开始“相当可用”。她把它称为 rule of thumb，而非普遍阈值；语言、方向、领域和错误代价都会改变可用性。}

\subsection{本章小结}

结果必须按 translation direction、resource level 与 metric 拆开。自动表格能支持快速系统比较，却不能证明用户感知质量、安全或所有方向一致提升；44\% 是有条件的相对改进，不是绝对质量承诺。

